Let $n\ge2$, $d=n-1$, $p_c=2n/d$, and fix $0<c_1<c_2$. On the spatial torus $\mathbb T^n=\mathbb R^n/(100\mathbb Z)^n$ write $X=(x,b)$ in an orthonormal frame with $x\in\mathbb R^d$. For each frequency with $\xi_b+|\xi|>0$ put $\kappa=(\xi_b+|\xi|)/\sqrt2$ and $v=\xi_x/\kappa$. There are constants $\gamma>0$, $K\ge1$, and $c,C>0$, depending only on $n,c_1,c_2$, with the following property. Let $0<h<1/(4K)$, $a=Kh$, and let $f$ be a trigonometric polynomial of $L^2$ norm one whose Fourier support consists of frequencies with $c_1\le h|\xi|\le c_2$, $\xi_b+|\xi|>0$, and $|v|\le\gamma$. Write $g_T=e^{-\mathrm i T|D|}f$.

Here is the energy used in the conclusion. For $z=(T,\tau)$, where $\tau$ lies in the quaternionic complex plane $\mathbb C_{\mathbf k}$ and $s=|\tau|\in(a,4)$, the auxiliary spatial coordinates are $(P,Z)\in\mathbb C_{\mathbf k}^{d}\times\mathbb C_{\mathbf k}$. Set
\[
Y=(\tau+Z)/\sqrt2,\quad S_1=(Z-\tau)/\sqrt2,\quad
y=|Y|,\quad W=P/Y,\quad \Phi=S_1+P\cdot P/(2Y),\quad Q=\Phi/a .
\]
Dots between complex vectors are bilinear. Take $w_0=1/10$ and let $\Omega_z$ be the torus times the auxiliary domain $Y\ne0$, $r=|W|^2/w_0^2+100|Q|^2-1<0$. Write $\mathrm d\zeta=\mathrm dX\,\mathrm dP\,\mathrm dZ$. The full quaternionic Hessian and gradient $H,G$ use the three real slots $(1,\mathbf i,\mathbf j)$ for each index, where $\mathbf k\mathbf i=\mathbf j$. In these slots with basis $e_\alpha$, their entries are
\[
H_{ab}(w)=\sum_{\alpha,\beta=0}^2 e_\alpha\overline{e_\beta}\,\partial_{a\alpha}\partial_{b\beta}w,\qquad
G_a(w)=\sum_{\alpha=0}^2e_\alpha\partial_{a\alpha}w .
\]
A subscript $\mathrm{fib}$ takes the spatial block at fixed $z$. Determinants and cofactors of these Hermitian matrices are the Moore determinant and its matrix derivative, respectively.

For $\ell>0$ let $J_\ell$ be the periodization of $c_n\ell^{-n}(1+|X/\ell|^2)^{-n-2}$, with $c_n$ chosen so the kernel integrates to one. Let $0\le g\le1$ be a fixed smooth cutoff equal to one on $|Q|\le1/100$ and zero on $|Q|\ge1/20$. The wave current $A^0=\chi\chi^*$ in null coordinates is given by
\[
\chi=\sum_\xi
\begin{pmatrix}1\\(v+W)/2\\-W\cdot v/2\end{pmatrix}
\iota_{Y/y}(\overline{\widehat f(\xi)e^{\mathrm i(X\cdot\xi-T|\xi|)}})
(\kappa/y)^{d/2}e^{-\kappa y|W-v|^2/2},
\qquad
\iota_e(\alpha+\mathrm i\beta)=\alpha+e\mathbf i\beta .
\]
Return to physical coordinates by $T=(t-s_0)/\sqrt2$, $b=(t+s_0)/\sqrt2$ in the coordinate index. With convolution in $X$, put
\[
A_{\mathrm w}=g(Q)(J_h*A^0),\qquad
\mu=(A_{\mathrm w})_{TT}+h^{50n}\operatorname{Cof}(H(r))_{TT},\qquad
B_0=\lambda|\Phi|^2,\quad \lambda=h^{-30n},\quad
D_w=H_{\mathrm{fib}}(B_0+w).
\]
For the smooth zero-boundary solution $u$ of $D_u>0$, $\det D_u=\lambda\mu$ on the fiber, set
\[
\mathcal F(T,\tau)=\int_{\Omega_z}(-u)\mu\,\mathrm d\zeta+
\frac1\lambda\int_{\Omega_z}u\left(\int_0^1\det D_{\theta u}\,\mathrm d\theta\right)\mathrm d\zeta .
\]
The energy $\mathcal F$ is nonnegative and smooth on the open parameter domain $\mathbb R_T\times\{\tau:a<|\tau|<4\}$. With all norms on the torus,
\begin{align*}
\mathcal F(0,\tau)&\le Ca^2\|f\|_{p_c}^{p_c} && (|\tau|=4a),\\
\mathcal F(T,\tau)&\ge ca^2\|g_T\|_{p_c}^{p_c} && (|\tau|=2a),\\
\Delta_z\mathcal F(T,\tau)&\le Cs^{-2}e^{-cs/h}\mathcal F(T,\tau) && (2a\le s\le3).
\end{align*}
The last two estimates hold for all real $T$, and $\Delta_z$ is the Euclidean Laplacian in $T$ and the two real coordinates of $\tau$.
